\begin{Solbook}{1}
 \par \noindent  \begin {listEX}[2] \item Ta có: \begin {align*} 12x-6&>26x+2024\\ 12x-26x&>2024+6\\ -14x&>2030\\ x&<-145 \end {align*} Vậy $x<-145$ \item $\begin {cases} 9x + 7y = 24 \\ 3x + 4y = -12 \end {cases}\\ \begin {cases} 9x + 7y = 24 \\ 9x + 12y = -36 \end {cases}\\ \begin {cases} 5y &= -60 \\ 3x &= -12-4y \end {cases}\\ \begin {cases} y &= -12 \\ x &= 12 \end {cases}\\ $ Vậy hệ phương trình có nghiệm duy nhất $(12;-12)$ \end {listEX} \par  \hfill \qedEX 
\end{Solbook}
\begin{Solbook}{2}
 \par \noindent  \begin {listEX}[2] \item \( \begin {aligned}[t] A &= 2\sqrt {20} - 3\sqrt {45} + \dfrac {1}{5}\sqrt {125}\\ &=4\sqrt {5}-9\sqrt {5}+\sqrt {5}=-4\sqrt {5} \end {aligned} \) \item \( \begin {aligned}[t] B &= \sqrt {(\sqrt {3}-2)^2}+ \dfrac {2\sqrt {3}}{\sqrt {3} + 1} + \dfrac {6}{\sqrt {3}}\\ &=|\sqrt {3}-2|+\dfrac {2\sqrt {3}(\sqrt {3}-1)}{3-1}+2\sqrt {3}\\ &=2-\sqrt {3}+\sqrt {3}(\sqrt {3}-1)+2\sqrt {3}\\ &=2-\sqrt {3}+3-\sqrt {3}+2\sqrt {3}\\ &=5 \end {aligned} \) \end {listEX}  \hfill \qedEX 
\end{Solbook}
\begin{Solbook}{3}
 \par \noindent  Gọi $x, y$ (học sinh) lần lượt là số học sinh nữ và học sinh nam lớp 9T $(x,y \in \mathbb {N^*})$\\ Số học sinh nam nhiều hơn số học sinh nữ là 8 học sinh: $-x+y=8 \quad (1)$\\ Tổng chiều cao của nữ lớp 9T: $161x$ (cm)\\ Tổng chiều cao của nam lớp 9T: $164y$ (cm)\\ Tổng chiều cao của cả lớp 9T: $162{,}8(x+y)=162{,}8x+168{,}2y$ (cm)\\ Suy ra: $161x+164y=162{,}8x+168{,}2y$ hay $-1{,}8x+1{,}2y=0 \quad (2)$\\ Từ (1) và (2), ta có hệ phương trình: $\begin {cases} -x+y&=8\\ -1{,}8x+1{,}2y&=0 \end {cases}$ \\ Giải hệ phương trình ta được: $x=16$, $y=24$ (thỏa mãn)\\ Vậy lớp 9T có 16 học sinh nữ. \par  \hfill \qedEX 
\end{Solbook}
\begin{Solbook}{4}
 \par \noindent  \begin {enumerate} \item Số tiền nhà trường phải trả theo phương án 1: $60x$ (đồng)\\ Số tiền nhà trường phải trả theo phương án 2: $500+54x$ (đồng) \item Để phương án 2 tiết kiệm hơn phương án 1, ta có bất phương trình: \begin {align*} 60x&>500+54x\\ 60x-54x&>500\\ 6x&>500\\ x&>\dfrac {250}{3} \end {align*} Vậy có ít nhất 84 học sinh tham gia thì nhà trường chọn phương án 2 sẽ tiết kiệm hơn. \end {enumerate}  \hfill \qedEX 
\end{Solbook}
\begin{Solbook}{5}
 \par \noindent  \begin {enumerate} \item Diện tích hình quạt $ABD$ được giới hạn bởi hai bán kính $AB$, $AD$ và cung nhỏ $BD$: \[\dfrac {\pi .3^2.60 ^{\circ }}{360 ^{\circ }}=\dfrac {3}{2}\pi \,(\text {cm}^2)\] \item Gọi $O$ là giao điểm $AC$ và $BD$. Khi đó $AC\perp BD$ tại $O$\\ $AC$ là đường phân giác của $\widehat {BAD}$ ($ABCD$ là hình thoi) suy ra $\widehat {BAO}=30^{\circ }$\\ Xét $\triangle ABO$ vuông tại $O$: $AO=AB\cdot \cos 30^{\circ }=\dfrac {3\sqrt {3}}{2}\,(\text {cm})$ và $BO=AB\cdot \sin 30^{\circ }=\dfrac {3}{2}\,(\text {cm})$\\ Suy ra $AC=3\sqrt {3}\,(\text {cm})$ và $BD=3\,(\text {cm})$\\ Diện tích hình thoi $ABCD$: $\dfrac {1}{2}AC\cdot BD=\dfrac {1}{2}\cdot 3\sqrt {3}\cdot 3=\dfrac {9\sqrt {3}}{2} \,(\text {cm}^2)$\\ Diện tích phần hình được tô màu: $\dfrac {9\sqrt {3}}{2}-\dfrac {3}{2}\pi \approx 3{,}1 \,(\text {cm}^2)$ \end {enumerate}  \hfill \qedEX 
\end{Solbook}
\begin{Solbook}{6}
 \par \noindent  \begin {center} \begin {tikzpicture}[scale=1, font=\footnotesize , >=stealth] \coordinate (O) at (0,0); \def \R {2.5} \draw [thick, black] (O) circle(\R ); \coordinate (A) at (-7.5,0); \path [name path=circleO] (O) circle(\R ); \path [name path=circleA] (-3.75,0) circle(3.75); \path [name intersections={of=circleO and circleA, by={B,C}}]; \coordinate (D) at ($-1*(B)$); \coordinate (E) at (-2.5,0); \coordinate (F) at (2.5,0); \path [name path=oa] (O)--(A); \path [name path=bc] (B)--(C); \path [name intersections={of=oa and bc, by=H}]; \path [name path=ab] (A)--(B); \path [name path=ed] (E)--(D); \path [name path=edline] (D)--($(D)!2!(E)$); \path [name intersections={of=ed and bc, by=K}]; \path [name intersections={of=edline and ab, by=M}]; \path [name path=mf] (M)--(F); \path [name path=be] (B)--(E); \path [name intersections={of=mf and be, by=N}]; \draw (A)--(B)--(C)--cycle; \draw (A)--(F)--(M); \draw (E)--(B)--(F); \draw (B)--(D)--(C); \draw (F)--(D)--(M); \draw (N)--(H); \foreach \p /\pos in {A/left,B/above,C/below,O/below right,D/below,E/below left,F/right,H/below right,K/below left,M/above,N/above} \filldraw [fill=white] (\p ) circle(1.5pt) node[\pos ] {$\p $}; \par \end {tikzpicture} \end {center} \begin {enumerate} \item Do $\begin {cases} AB=AC \,(\text {Tính chất hai tiếp tuyến cắt nhau})\\ OB=OC \, (\text {Bán kính đường tròn}) \end {cases}$ \par Nên $OA$ là đường trung trực của $BC$ \par Suy ra $ OA \perp BC$ tại $H$. \par Ta có: $\widehat {BED}=90^{\circ }$ (chắn nửa đường tròn) \par Xét $\triangle BED$ có $\widehat {BED}=90^{\circ }$, suy ra $\triangle BED$ vuông tại $E$. \item Xét $\triangle MBE$ vuông tại $E$ và $\triangle MDB$ vuông tại $B$ có $\widehat {BMD}$ là góc chung \par Suy ra $\triangle MBE \backsim \triangle MDB$ (g.g) \par Suy ra $\dfrac {MB}{MD}=\dfrac {ME}{MB}$ suy ra $MB^2=MD.ME$ \par Do $\widehat {ABE}=\widehat {EBD}$ ($\triangle MBE \backsim \triangle MDB$) và $\widehat {EBK}=\widehat {EBD}$ (cùng chắn $ \stackrel \frown {EC}$) \par Nên $\widehat {ABE}=\widehat {EBK}$ \item $\triangle MBK$ có $BE$ vừa là đường cao, vừa là đường phân giác nên $\triangle MBK$ cân tại $B$ \par Suy ra $BE$ là trung tuyến của $\triangle MBK$ do đó $E$ là trung điểm $KM$ suy ra $KE=ME$ \par Ta có: $\widehat {BCD}=90^{\circ }, \widehat {FDE}=90^{\circ }$ (chắn nửa đường tròn) \par Xét tứ giác $BFDE$ có $\widehat {BCD}=\widehat {FDE}=\widehat {BED}=90^{\circ }$ suy ra $BFDE$ là hình chữ nhật \par Do $\begin {cases} OA \perp BC \,(\text {cmt})\\ DC \perp BC \,(\widehat {BCD}=90^{\circ }) \end {cases}$ nên $OA \parallel DC $ \par $HE \parallel DC\, (OA \parallel BC)$ suy ra $\dfrac {DE}{EK}=\dfrac {HC}{HK}$ (Định lý Thales) \par \par Mà $ME=EK$ và $HC=HB$ suy ra $\dfrac {DE}{ME}=\dfrac {HB}{HK}$\\ \par $\dfrac {DE}{ME}+1=\dfrac {HB}{HK}+1$\\ \par $\dfrac {DE+ME}{ME}=\dfrac {HB+HK}{HK}$\\ \par $\dfrac {MD}{ME}=\dfrac {BK}{HK} \quad (1)$ \par $BE \perp DE\, (\widehat {BED}=90^{\circ })$ suy ra $BE\parallel FD$ hay $NE\parallel FD$ suy ra $\dfrac {FD}{NE}=\dfrac {MD}{ME}$ \par Mà $FD=BE$ ($BFDE$ là hình chữ nhật) suy ra $\dfrac {BE}{NE}=\dfrac {MD}{ME} \quad (2)$ \par Từ (1) và (2), suy ra: $\dfrac {BK}{HK}=\dfrac {BE}{NE}$ \par Áp dụng định lý Thales đảo, ta có: $NH\parallel KE$ hay $NH \parallel MK$ \par \end {enumerate}  \hfill \qedEX 
\end{Solbook}
